% discussion.tex — Paper 6, Section VII: Discussion
% ASSERT_CONVENTION: natural_units=natural, metric_signature=mostly_plus, entropy_base=nats, modular_hamiltonian=K_minus_ln_rho, coupling_convention=H_sum_hxy, state_normalization=Tr_rho_1
\section{Discussion}
\label{sec:discussion}

\subsection{What is derived versus what is input}
\label{sec:derived-vs-input}

The derivation chain summarized in Table~\ref{tab:chain} establishes
Einstein's field equations from self-modeling through eight links.  It is
important to state precisely what is derived, what is input, and what is
standard methodology.

\emph{Derived from self-modeling.}---The local algebra
$M_n(\mathbb{C})^{sa}$ with the L\"uders sequential product is forced by
faithful self-modeling~\cite{Paper5}. Diagonal $U(n)$ covariance of the
boundary coupling is derived from the single-probe structure of the
self-model (Sec.~\ref{sec:hamiltonian}).  The SWAP interaction
Hamiltonian $H = \sum_{\langle x,y\rangle} J\,\swapop$ then follows from
Schur--Weyl duality.  The area-law entanglement structure follows from
known properties of the Heisenberg model: the Calabrese--Cardy formula
and WVCH mutual information bound~\cite{WVCH2008}, combined with the
exact entanglement first law
$\delta\ent = \delta\langle\modH\rangle$
(Sec.~\ref{sec:arealaw}).

\emph{Input: lattice topology.}---The graph $G = (V, E)$ defining which
sites are nearest neighbors is not derived from self-modeling.  The paper
derives the emergent \emph{metric}---the assignment of distances
compatible with the entanglement structure---but not the
\emph{topology}.  The spatial dimension is set by the graph structure:
a $d$-dimensional regular lattice yields $d$ emergent spatial dimensions.
This is an irreducible input to the framework.

\emph{Input: antiferromagnetic coupling.}---The sign of~$J$ is not
determined by self-modeling.  Only $J > 0$ (antiferromagnetic) supports
nontrivial entanglement; the ferromagnetic ground state is a product
state with zero entanglement (Sec.~\ref{sec:lattice}).  We work with
$J > 0$ throughout; this is an input.

\emph{Standard methodology: Wilsonian continuum limit.}---The passage
from a lattice theory to a smooth Riemannian manifold at long wavelengths
$L \gg a$ (where $a$ is the lattice spacing) employs a Wilsonian
universality argument.  This is standard methodology in lattice QCD,
causal dynamical triangulations, and Regge calculus.  However, its
applicability to the self-modeling Hamiltonian in $d \geq 2$ is not
guaranteed---see Gap~1 below.

\emph{Structural identification: MVEH via thermal time.}---The maximal
vacuum entanglement hypothesis does not appear as a numbered assumption
in our framework.  In the absence of pre-existing geometry, the modular
flow of the self-modeling lattice state defines the emergent geometric
flow, following the thermal time hypothesis~\cite{ConnesRovelli1994}.
This is the only available identification of geometry with state in a
pre-geometric framework, and the state so identified necessarily
satisfies MVEH.  We regard this as a structural identification---necessary
but not mathematically derived.  The existence and uniqueness of a
geometry-defining state are open questions.
Section~\ref{sec:equilibrium} discusses this identification in detail,
including the Sorce caveat~\cite{Sorce2024} on conformal symmetry
requirements for geometric modular flow.

\subsection{Remaining gaps}
\label{sec:gaps}

Three gaps remain in the derivation chain.

\paragraph{Gap 1: Continuum limit and $d \geq 2$ infrared behavior
(principal open problem).}
The Wilsonian argument that the lattice produces a smooth Riemannian
manifold $(M, \gab)$ at long wavelengths is the \emph{principal open
problem} of this paper.  Both Route~A and Route~B require a smooth
emergent geometry; neither derives it.  A rigorous proof would require
showing that lattice correlation functions converge to those of a local
quantum field theory on a Riemannian manifold as $a \to 0$.

In $d \geq 2$, this is compounded by the infrared behavior of the
Heisenberg antiferromagnet: the $d \geq 2$ model is expected to exhibit
N\'eel order (spontaneous symmetry breaking of $SU(n)$) rather than
conformal symmetry.  The resulting effective theory is a nonlinear sigma
model with a mass gap, not a CFT.  Whether this produces a smooth
Riemannian geometry at long wavelengths, and if so with what
signatures---a mass gap may correspond to a cosmological constant, for
instance---is an open question in condensed matter physics that we
cannot resolve here.

This is the same wall that the entire ``gravity from entanglement''
program hits.  Jacobson's argument~\cite{Jacobson2016} assumes smooth
geometry.  The holographic results
of~\cite{LMVR2014, Faulkner2014} assume AdS/CFT.  The CCM
reconstruction~\cite{CCM2017} assumes a factored Hilbert space with
specific entanglement structure.  No program in this field derives the
continuum limit from a UV theory.  Our contribution is the UV
completion (L1--L5), not the continuum limit.  The lattice provides the
UV; the continuum limit is the interface to GR; and that interface is
open for everyone.

\paragraph{Gap 2: Conformal approximation (Route A only).}
The Casini--Huerta--Myers conformal modular
Hamiltonian~\cite{CasiniHuertaMyers2011}, used in Route~A
(Sec.~\ref{sec:route-a}) to evaluate the matter entropy variation, is
exact only for conformal field theories in causal diamonds.  For
nonconformal fields, corrections of order
$O\!\left((m\ell)^{2\Delta}\right)$ appear~\cite{Speranza2016}, where
$m$ is the mass scale breaking conformal invariance and $\ell$ is the
ball radius.  In $d = 1$, the $SU(2)_1$ WZW CFT makes this gap vanish.
In $d \geq 2$, the likely N\'eel order makes the conformal approximation
problematic.  This gap is specific to Route~A; Route~B
(Sec.~\ref{sec:route-b}) circumvents it via Lovelock uniqueness.

\paragraph{Gap 3: Tensoriality assumption (Route B only).}
The Lovelock route assumes that the entanglement-geometry equation is a
symmetric 2-tensor equation with at most second derivatives
(Sec.~\ref{sec:route-b}).  Symmetry follows from the metric; the
second-derivative restriction follows from the entanglement first law
being first-order in the perturbation, coupling to the modular
Hamiltonian (local at leading order in the derivative expansion).
Divergence-freeness is then a consequence of Lovelock's theorem
itself (the contracted Bianchi identity), not an additional input.
These justifications are physically reasonable but constitute an
assumption, not a derivation from self-modeling.  If the
entanglement-geometry relation involves higher derivatives or non-local
terms, Lovelock's theorem does not apply and the resulting equation
could include higher Lovelock terms (in $d + 1 > 4$) or other
modifications.  This gap is specific to Route~B; Route~A derives the
equation form explicitly.

\paragraph{Note on MVEH.}
MVEH is not listed as a gap; it is reframed as a structural
identification via the Connes--Rovelli thermal time
hypothesis~\cite{ConnesRovelli1994}
(Sec.~\ref{sec:equilibrium}).  This identification is necessary
in a pre-geometric framework---there is no other candidate for
the vacuum---but it is not a mathematical derivation: the existence
and uniqueness of a geometry-defining state are open questions.
Numerical evidence from Sec.~\ref{sec:numerical} is consistent
with MVEH: all Hamiltonian perturbations tested decrease $\ent(A)$.

\subsection{Comparison with related work}
\label{sec:comparison}

\paragraph{Jacobson (2016)~\cite{Jacobson2016}.}
Jacobson's entanglement equilibrium derivation assumes three inputs:
area-law entanglement entropy, the entanglement first law, and MVEH.
Our contribution is the UV completion upstream of these inputs: we
\emph{derive} the area-law entanglement structure from the self-modeling
lattice (Links~L1--L5 in Table~\ref{tab:chain}), rather than assuming it.
Jacobson's derivation was UV-agnostic---it applies to any UV completion
satisfying its inputs.  We provide a specific, forced UV completion.
Route~A follows Jacobson's argument directly; Route~B provides an
independent uniqueness argument via Lovelock's theorem that does not
require the conformal modular Hamiltonian.

\paragraph{Cao--Carroll--Michalakis (2017)~\cite{CCM2017}.}
CCM derive the spatial constraint equation (the Hamiltonian constraint of
general relativity) from the entanglement structure of a
finite-dimensional Hilbert space.  Our result gives the full
\emph{spacetime} Einstein equation, including the dynamical sector, not
only the spatial constraint.  The spatial part of our Einstein equation is
consistent with CCM's constraint.  Both approaches share the insight that
entanglement structure encodes geometry; the key difference is that CCM
reconstruct spatial geometry from entanglement, while we (following
Jacobson) derive the full gravitational dynamics from entanglement
equilibrium.

\paragraph{Lashkari--McDermott--Van Raamsdonk (2014)~\cite{LMVR2014}.}
LMVR showed that the entanglement first law in holographic conformal field
theories implies the \emph{linearized} Einstein equation.  Our
derivation, following Jacobson~\cite{Jacobson2016}, extends this to the
full \emph{nonlinear} Einstein equation via the entanglement equilibrium
condition.  The LMVR result is recovered as the linearized limit of our
Eq.~(\ref{eq:einstein}) around flat space.  Crucially, the LMVR derivation
assumes AdS/CFT, while ours is non-holographic.

\paragraph{Faulkner \textit{et~al.}\ (2014)~\cite{Faulkner2014}.}
Faulkner~\textit{et~al.}\ derived the full nonlinear Einstein equations
from entanglement in the holographic setting, extending LMVR beyond the
linearized regime.  Their derivation, like LMVR's, assumes the AdS/CFT
correspondence.  Our approach is non-holographic: the self-modeling
lattice does not assume a bulk/boundary duality.  The self-modeling
framework provides a UV completion in its own right, independent of
holography.

\subsection{Connection to Paper~5}
\label{sec:paper5-connection}

Paper~5~\cite{Paper5} showed that faithful self-modeling forces quantum
mechanics: the state space is $M_n(\mathbb{C})^{sa}$ with the L\"uders
sequential product and the conjugate-transpose involution.  This paper
shows that self-modeling forces general relativity: the locality of
self-modeling produces area-law entanglement, and either Jacobson's
argument (Route~A, conformal regime) or Lovelock uniqueness (Route~B,
$d \geq 2$) converts this to Einstein's equations.

Together, the two papers establish that the self-modeling
principle---combined with a lattice topology---yields both quantum
mechanics and general relativity as necessary consequences.  Quantum
mechanics emerges from the algebraic structure of single-site
self-modeling (the $B$--$M$ boundary interaction).  General relativity
emerges from the entanglement structure of the multi-site lattice (the
inter-site boundary interactions).

\subsection{What this paper does not claim}
\label{sec:not-claimed}

We state explicitly what lies outside the scope of this paper:
\begin{itemize}
\item \textbf{Newton's constant.}  The value $\Gnewton = 1/(4\eta)$
  depends on the entanglement entropy density~$\eta$, which in turn
  depends on the lattice local dimension~$n$ and the lattice spacing~$a$.
  These are not determined by self-modeling.
\item \textbf{Spacetime dimension.}  The spatial dimension~$d$ is set by
  the lattice topology $G = (V, E)$, which is an irreducible input.
  Self-modeling does not select $d = 3$.
\item \textbf{Cosmological constant.}  $\Lambda$ appears as an
  undetermined integration constant in the Jacobson
  derivation~\cite{Jacobson2016, JacobsonSperanza2019}.  We do not
  predict its value.
\item \textbf{Null energy condition.}  The sign chain producing
  attractive gravity (Sec.~\ref{sec:einstein}) assumes the NEC
  ($\Tab\, k^a k^b \geq 0$ for null $k^a$).  The tensor
  equation~\eqref{eq:einstein} holds without the NEC; only the sign
  interpretation requires it.
\item \textbf{Proof of MVEH.}  We do not prove the maximal vacuum
  entanglement hypothesis from self-modeling.  We identify it as the
  only available vacuum selection criterion in a pre-geometric framework
  via the thermal time hypothesis.  This is a structural identification,
  not a derivation.  The existence and uniqueness of a geometry-defining
  state are open questions.
\end{itemize}

\subsection{Future work}
\label{sec:future}

Several directions are natural extensions of this work.

\emph{Spatial dimension from self-modeling.}---Can constraints on the
lattice topology $G = (V, E)$ be derived from self-modeling requirements?
If so, this would reduce the input set further, potentially selecting
$d = 3$ spatial dimensions.

\emph{Matter content.}---The self-modeling lattice produces a gravitational
sector.  An important question is whether constraints on matter content
(gauge groups, fermion representations) also follow from self-modeling,
analogous to anomaly cancellation constraints in string theory.

\emph{Spectral gap.}---Proving that the self-modeling Hamiltonian has a
spectral gap above the ground state (for appropriate sign of~$J$ and
appropriate dimension) would strengthen the area-law argument, bringing
the Hastings theorem~\cite{Hastings2007} and the Brand\~ao--Horodecki
result~\cite{BrandaoHorodecki2013} to bear directly.

\emph{Larger numerical simulations.}---The exact diagonalization results
in Sec.~\ref{sec:numerical} are limited to $N \leq 20$ sites.  Tensor
network methods (DMRG, TEBD) could push to $N \sim 10^2$--$10^3$ in one
dimension, providing stronger evidence for the area-law scaling and MVEH.

\emph{Thermodynamic arrow of time.}---The thermal time hypothesis
identifies physical time with modular flow.  The connection between this
identification and the thermodynamic arrow of time in the self-modeling
framework is an open and natural question.
